\section{Desarrollo Frontend y Experiencia de Usuario (UX)}
\label{sec:frontend_sprint3}

La capa de presentación fue expandida para orquestar la carga dinámica de evidencias y la visualización del portafolio público, implementando patrones de UX defensiva para prevenir bloqueos de renderizado.

\subsection{Gestión Interactiva de Proyectos (GAN-001, GAN-002)}
Se desarrollaron formularios controlados para la creación y edición de proyectos, integrados directamente con el estado global de la aplicación. Para la gestión de evidencias, se implementó una interfaz de \textit{Drag \& Drop} que permite al perfil reordenar visualmente sus adjuntos y enlaces, garantizando una presentación coherente.

\begin{AlertaDefecto}
Durante la integración de Iframes para embeber videos y prototipos (ej. YouTube, Figma), se detectó que enlaces mal formados rompían el layout principal. Se mitigó envolviendo el componente en un \textit{ErrorBoundary} e implementando un \textit{fallback} UI ('Contenido no disponible') junto con el uso masivo de \textit{Optional Chaining} (\comando{?.}) en los flujos de datos para evitar \textit{crashes} si faltan propiedades nulas (GAN-006).
\end{AlertaDefecto}

\subsection{Optimización de Carga con Dropzone (GAN-005)}
Para la adjunción de archivos físicos, se construyó un componente \comando{Dropzone} altamente intuitivo. Para reducir la latencia y el consumo innecesario de ancho de banda, el componente ejecuta una validación local mediante JavaScript (peso máximo y extensión permitida) previa a la solicitud HTTP. Una barra de progreso reactiva proporciona feedback continuo durante la transferencia de red.

\subsection{Visor de Catálogo y Seguridad de Sesión (GAN-003, GAN-004)}
La experiencia del reclutador fue materializada en una grilla responsiva basada en tarjetas (Cards), optimizando el \textit{viewport} al ocultar la barra lateral en resoluciones de escritorio. En paralelo, se implementó un panel de administración (\comando{AdminDomainsPage}) con un diseño OLED oscuro para la gestión de dominios permitidos.

Como medida de seguridad proactiva, se inyectó un \textit{Idle Timer} en el \textit{layout} autenticado. Tras 15 minutos de inactividad, la sesión del perfil se invalida localmente, purgando el estado global y forzando una redirección al \textit{login}.

\begin{TablaBacklog}
    GAN-001 & Frontend: Maquetar formulario de proyecto con estado global. & Alt & 1.5 & 3 & \estadoPass \\ \hline
    GAN-003 & Frontend: Construir AdminDomainsPage con diseño OLED. & Alt & 1.5 & 3 & \estadoPass \\ \hline
    GAN-005 & Frontend: Dropzone intuitivo y validación local rápida. & Alt & 2.0 & 3 & \estadoPass \\ \hline
    GAN-006 & Frontend: Optional Chaining en UI para evitar crashes. & Med & 1.0 & 3 & \estadoPass \\ \hline
\end{TablaBacklog}